% GENERATED FILE. Edit canonical lessons, then run npm run generate:books.
% canonical-source-hash: fnv1a64:1c7842fb6bb1f1d6
% canonical-lessons: UR-C352-mazmun, UR-C352-maqala, UR-C352-danishvar, UR-C352-ilm, UR-C352-chhantna

\chapter{Ask an Expert}
\label{ch:ur-ask-an-expert}
\hlchaptermodality{352}

\begin{chapteropening}
\textbf{By the end of this chapter:} \emph{I can say who knows best and where they said it.}

Complete the last lesson of chapter 352: I can say who knows best and where they said it.
\end{chapteropening}

\section[\texorpdfstring{\ur{مضمون} (mazmūn)}{مضمون (mazmūn)}]{\ur{مضمون} (mazmūn) — an article; an essay}

\label{lesson:UR-C352-mazmun}



\begin{quote}
Before the new one: say the Urdu for to compare, then the Urdu for to assess.
\end{quote}

\subsection*{You'll want to know: \ur{مضمون}}
\textbf{\ur{مضمون}} — \emph{mazmūn} — \textquotedblleft{}an article; an essay\textquotedblright{}.

\begin{quote}
\ur{میں} \ur{نے} \ur{اخبار} \ur{میں} \ur{اس} \ur{کا} \ur{مضمون} \ur{پڑھا۔}
\end{quote}

\begin{quote}
\emph{maiṅ ne akhbār meṅ us kā mazmūn paṛhā.}
\end{quote}

\begin{quote}
I read her article in the newspaper.
\end{quote}

Masculine.

\begin{cousinweb}
\textbf{\ur{مضمون}} (\emph{mazmūn}) is usually traced, through Persian, to Arabic \emph{maḍmūn}, \textquotedblleft{}contents, what is contained\textquotedblright{}.
\end{cousinweb}

\begin{grammarlens}[title={What you've built}]
One more word for giving reasons and opinions.
\end{grammarlens}

\subsection*{Your turn}
\begin{itemize}
\raggedright
  \item \emph{Say it:} \emph{mazmūn}
  \item \emph{Say it:} \emph{mazmūn}, once more
  \item \emph{Say it:} \emph{maiṅ ne akhbār meṅ us kā mazmūn paṛhā.} — I read her article in the newspaper.
  \item \emph{Recall:} say the Urdu for to compare, then the Urdu for to assess, then say \emph{mazmūn} again
\end{itemize}

\begin{tcolorbox}[breakable,colback=teal!4,colframe=teal!35!black,title={Before you move on}]
What does \ur{مضمون} mean? (\textquotedblleft{}An article; an essay\textquotedblright{}.) Say it once more.
\end{tcolorbox}
\section[\texorpdfstring{\ur{مقالہ} (maqālā)}{مقالہ (maqālā)}]{\ur{مقالہ} (maqālā) — a thesis; a research paper}

\label{lesson:UR-C352-maqala}



\begin{quote}
Before the new one: say the Urdu for to assess, then the Urdu for an article.
\end{quote}

\subsection*{You'll want to know: \ur{مقالہ}}
\textbf{\ur{مقالہ}} — \emph{maqālā} — \textquotedblleft{}a thesis; a research paper\textquotedblright{}.

\begin{quote}
\ur{اس} \ur{نے} \ur{اپنا} \ur{مقالہ} \ur{ختم} \ur{کیا۔}
\end{quote}

\begin{quote}
\emph{us ne apnā maqālā khatm kiyā.}
\end{quote}

\begin{quote}
She finished her thesis.
\end{quote}

Masculine.

\begin{cousinweb}
\textbf{\ur{مقالہ}} (\emph{maqālā}) is usually traced, through Persian, to Arabic \emph{maqāla}, \textquotedblleft{}a saying; a treatise\textquotedblright{}.
\end{cousinweb}

\begin{grammarlens}[title={What you've built}]
One more word for giving reasons and opinions.
\end{grammarlens}

\subsection*{Your turn}
\begin{itemize}
\raggedright
  \item \emph{Say it:} \emph{maqālā}
  \item \emph{Say it:} \emph{maqālā}, once more
  \item \emph{Say it:} \emph{us ne apnā maqālā khatm kiyā.} — She finished her thesis.
  \item \emph{Recall:} say the Urdu for to assess, then the Urdu for an article, then say \emph{maqālā} again
\end{itemize}

\begin{tcolorbox}[breakable,colback=teal!4,colframe=teal!35!black,title={Before you move on}]
What does \ur{مقالہ} mean? (\textquotedblleft{}A thesis; a research paper\textquotedblright{}.) Say it once more.
\end{tcolorbox}
\section[\texorpdfstring{\ur{دانشور} (dānishvar)}{دانشور (dānishvar)}]{\ur{دانشور} (dānishvar) — an intellectual; a thinker}

\label{lesson:UR-C352-danishvar}



\begin{quote}
Before the new one: say the Urdu for an article, then the Urdu for a thesis.
\end{quote}

\subsection*{You'll want to know: \ur{دانشور}}
\textbf{\ur{دانشور}} — \emph{dānishvar} — \textquotedblleft{}an intellectual; a thinker\textquotedblright{}.

\begin{quote}
\ur{شہر} \ur{کے} \ur{دانشور} \ur{اس} \ur{بات} \ur{پر} \ur{متفق} \ur{ہیں۔}
\end{quote}

\begin{quote}
\emph{shahr ke dānishvar is bāt par muttafiq haiṅ.}
\end{quote}

\begin{quote}
The city's thinkers agree on this point.
\end{quote}

Masculine or feminine, by the person meant.

\begin{cousinweb}
\textbf{\ur{دانشور}} (\emph{dānishvar}) is usually traced to Persian \emph{dānishvar}, \textquotedblleft{}possessing knowledge\textquotedblright{}, built on \emph{dānish}, \textquotedblleft{}knowledge\textquotedblright{}.
\end{cousinweb}

\begin{grammarlens}[title={What you've built}]
One more word for giving reasons and opinions.
\end{grammarlens}

\subsection*{Your turn}
\begin{itemize}
\raggedright
  \item \emph{Say it:} \emph{dānishvar}
  \item \emph{Say it:} \emph{dānishvar}, once more
  \item \emph{Say it:} \emph{shahr ke dānishvar is bāt par muttafiq haiṅ.} — The city's thinkers agree on this point.
  \item \emph{Recall:} say the Urdu for an article, then the Urdu for a thesis, then say \emph{dānishvar} again
\end{itemize}

\begin{tcolorbox}[breakable,colback=teal!4,colframe=teal!35!black,title={Before you move on}]
What does \ur{دانشور} mean? (\textquotedblleft{}An intellectual; a thinker\textquotedblright{}.) Say it once more.
\end{tcolorbox}
\section[\texorpdfstring{\ur{علم} (ilm)}{علم (ilm)}]{\ur{علم} (ilm) — knowledge; learning}

\label{lesson:UR-C352-ilm}



\begin{quote}
Before the new one: say the Urdu for a thesis, then the Urdu for an intellectual.
\end{quote}

\subsection*{You'll want to know: \ur{علم}}
\textbf{\ur{علم}} — \emph{ilm} — \textquotedblleft{}knowledge; learning\textquotedblright{}.

\begin{quote}
\ur{کتابوں} \ur{سے} \ur{علم} \ur{ملتا} \ur{ہے۔}
\end{quote}

\begin{quote}
\emph{kitāboṅ se ilm miltā hai.}
\end{quote}

\begin{quote}
Books give us knowledge.
\end{quote}

Masculine.

\begin{cousinweb}
\textbf{\ur{علم}} (\emph{ilm}) is usually traced, through Persian, to Arabic \emph{ʿilm}, \textquotedblleft{}knowledge\textquotedblright{}.
\end{cousinweb}

\begin{grammarlens}[title={What you've built}]
One more word for giving reasons and opinions.
\end{grammarlens}

\subsection*{Your turn}
\begin{itemize}
\raggedright
  \item \emph{Say it:} \emph{ilm}
  \item \emph{Say it:} \emph{ilm}, once more
  \item \emph{Say it:} \emph{kitāboṅ se ilm miltā hai.} — Books give us knowledge.
  \item \emph{Recall:} say the Urdu for a thesis, then the Urdu for an intellectual, then say \emph{ilm} again
\end{itemize}

\begin{tcolorbox}[breakable,colback=teal!4,colframe=teal!35!black,title={Before you move on}]
What does \ur{علم} mean? (\textquotedblleft{}Knowledge; learning\textquotedblright{}.) Say it once more.
\end{tcolorbox}
\section[\texorpdfstring{\ur{چھانٹنا} (chhānṭnā)}{چھانٹنا (chhānṭnā)}]{\ur{چھانٹنا} (chhānṭnā) — to pick out; to sort (the good from the bad)}

\label{lesson:UR-C352-chhantna}



\begin{quote}
Before the new one: say the Urdu for an intellectual, then the Urdu for knowledge.
\end{quote}

\subsection*{You'll want to know: \ur{چھانٹنا}}
\textbf{\ur{چھانٹنا}} — \emph{chhānṭnā} — \textquotedblleft{}to pick out; to sort (the good from the bad)\textquotedblright{}.

\begin{quote}
\ur{ماہر} \ur{نے} \ur{اچھے} \ur{نمونے} \ur{چھانٹے۔}
\end{quote}

\begin{quote}
\emph{māhir ne achchhe namūne chhānṭe.}
\end{quote}

\begin{quote}
The expert picked out the good samples.
\end{quote}

Past, with \ur{نے} (he or she …): \ur{اس} \ur{نے} \ur{سیب} \ur{چھانٹے}, \emph{us ne seb chhānṭe}.

\begin{cousinweb}
\textbf{\ur{چھانٹنا}} (\emph{chhānṭnā}) is inherited: CDIAL lists the Hindi cousin \emph{chā̃ṭnā}, \textquotedblleft{}to sift, to select\textquotedblright{}, under \emph{*chaṇṭ}, a form it reconstructs (CDIAL 4970).
\end{cousinweb}

\begin{grammarlens}[title={What you've built}]
One more word for giving reasons and opinions.
\end{grammarlens}

\subsection*{Your turn}
\begin{itemize}
\raggedright
  \item \emph{Say it:} \emph{chhānṭnā}
  \item \emph{Say it:} \emph{chhānṭnā}, once more
  \item \emph{Say it:} \emph{māhir ne achchhe namūne chhānṭe.} — The expert picked out the good samples.
  \item \emph{Recall:} say the Urdu for an intellectual, then the Urdu for knowledge, then say \emph{chhānṭnā} again
\end{itemize}

\begin{tcolorbox}[breakable,colback=teal!4,colframe=teal!35!black,title={Before you move on}]
What does \ur{چھانٹنا} mean? (\textquotedblleft{}To pick out; to sort (the good from the bad)\textquotedblright{}.) Say it once more.
\end{tcolorbox}
